\documentclass[11pt]{article}
\usepackage{series}
\usepackage{booktabs,tabularx,array}
\usepackage{xcolor}

\newcommand{\MTT}{\mathrm{MTT}}
\newcommand{\Adm}{\mathcal A}
\newcommand{\Image}{\mathcal R}
\newcommand{\Motif}{\mathcal M}
\newcommand{\Graphs}{\mathcal G}
\newcommand{\Coherent}{\mathcal X}
\newcommand{\Effective}{\mathcal Y}
\newcommand{\Prob}{\mathbb P}
\newcommand{\E}{\mathbb E}
\newcommand{\R}{\mathbb R}
\newcommand{\Sone}{S^1}

\title{\textbf{Admissible Image Geometry for Network Motifs in Modal Triplet Theory}\\
\large Conditional Correlations, Sprout Bounds, and What Projection Alone Does Not Imply}
\author{Peter Nero}
\date{Version 2, July 2026}

\begin{document}
\maketitle

\begin{abstract}
Projection can discard information without selecting a network geometry.
This paper develops a rigorous Modal Triplet Theory (MTT) framework in which
observable motif constraints are determined from the actual image
\[
  \Image=F(\Adm)\subseteq\Motif,
\]
where $\Adm$ is a declared admissible source domain and
$F$ includes projection, network extraction, and motif measurement.
Compatibility constraints are properties of $\Image$, not consequences of
noninvertibility by itself.  We give an exact recombination test for support
nonfactorization, a graph-of-a-map sufficient condition, and a regular-value
description of local image geometry.  A covering-map counterexample proves
that failure of a continuous section is compatible with a full product
image.  We also separate three distinct objects: an embedded
physical network, a constrained motif support, and a statistical
conditional-dependence graph.

For a declared local reserve model, we derive a useful conditional result.
If the leading coherence deficit is
\[
  Q(\beta,q_\parallel,q_\perp)
  =a\beta^2+b(q_\parallel^2+q_\perp^2)+cq_\parallel^2,
  \qquad a,b,c>0,
\]
then every admissible branch of thickness
$\rho=(q_\parallel^2+q_\perp^2)^{1/2}$ satisfies
$\rho\leq\sqrt{C/b}$.  Near saturation, the dominant channel is forced to
remain nearly straight and the branch nearly orthogonal.  The estimate is
stable under a relative remainder bound.  This proves the square-root law
inside the stated quadratic model; it does not derive the model, its
coefficients, a branch-creation law, or universal empirical prevalence.
Those require a selected network source, Hessian, extraction pipeline, and
probability or growth dynamics.  The result is a falsifiable conditional
encoding rather than a universal theory of network morphology.
\end{abstract}

\section*{Version 2 revision note}

\begin{description}[style=nextline]
\item[Supersedes.]
Version 1, DOI
\url{https://doi.org/10.5281/zenodo.18274629}.
\item[Reason.]
Version 1 treated the absence of a global reconstruction map as if it forced
a proper, nonfactorizing motif image and then used that inference to derive
specific graph geometry, sprout prevalence, and a universal square-root
law.  None of those implications follows from noninvertibility alone.
\item[Resolution.]
Version 2 defines the image $\Image=F(\Adm)$ explicitly, supplies direct
tests for its geometry, separates image constraints from graph extraction
and statistical dependence, and states the sprout result as a theorem of a
declared positive quadratic reserve model.
\item[Retained content.]
The useful local idea survives: a selected positive Hessian can make bending
and longitudinal branch growth more expensive than a transverse thin mode.
Near a small reserve boundary, that anisotropy yields a controlled,
approximately orthogonal sprout bound.
\item[Open boundary.]
MTT has not yet selected a physical network source, the map $F$, the
coefficients $a,b,c$, a capacity calibration, or a growth law for any
biological or engineered system.  Cross-domain universality and empirical
prevalence are therefore open.
\end{description}

\tableofcontents

\section{The question this paper can answer}

Networks appear in vasculature, leaf venation, fungal growth, neural
arborization, infrastructure, and many other settings.  Their common visual
vocabulary does not imply a common microscopic cause.  Even within transport
networks, optimization, adaptation, damage tolerance, fluctuating loads, and
tissue growth can produce different structures
\cite{katifori2010,ronellenfitsch2016,corson2010}.  Network motifs are useful
empirical summaries, but a repeated motif is not itself a derivation of its
mechanism \cite{milo2002,newman2018}.

The narrow MTT question is this:

\begin{quote}
Given a selected source domain, a selected projection and extraction map,
and a certified local reserve function, what restrictions follow for the
observable motif image?
\end{quote}

This question is mathematically meaningful.  It is also weaker than the
claim that projection alone explains real network morphology.  The
difference matters.  A many-to-one map tells us that some source information
was forgotten.  It does not tell us which motifs remain, which edges appear,
or how probability is distributed over the surviving motifs.

Version 2 therefore uses four rules.

\begin{enumerate}
\item The realizable set is computed from the image of a declared map.
\item Nonfactorization is proved by a failed recombination or an explicit
coupled constraint.
\item Graph topology is supplied by an extraction rule, not inferred from
correlation.
\item A prevalence statement requires a measure or dynamics in addition to
a feasible set.
\end{enumerate}

These rules turn the paper from a universality claim into a reusable
conditional framework.

\section{Typed source, image, and graph objects}

\subsection{The full map}

Let $\Coherent$ be a source configuration space and let
$\Adm\subseteq\Coherent$ be the domain on which the relevant MTT source,
projector, and stability conditions have been verified.  Let
\[
  P:\Adm\longrightarrow\Effective
\]
be the effective projection.  Network extraction is a separate map
\[
  \Gamma:\Effective\longrightarrow\Graphs,
\]
where an element of $\Graphs$ contains a graph, an embedding when relevant,
and edge attributes such as radius or conductivity.  Finally, let
\[
  m:\Graphs\longrightarrow\Motif
\]
be the declared motif-measurement map.  The complete observable map is
\[
  F=m\circ\Gamma\circ P:\Adm\longrightarrow\Motif.
\]

\begin{definition}[Admissible motif image]
The realizable motif support of the declared construction is
\[
  \Image:=F(\Adm)\subseteq\Motif.
\]
If the source carries a probability measure $\nu$, its observable law is
$\mu=F_\ast\nu$.  The pair $(\Image,\mu)$ contains more information than
either object alone.
\end{definition}

Every arrow in this definition has an owner.  Changing the skeletonization
algorithm changes $\Gamma$.  Changing how a junction is measured changes
$m$.  Changing the coherent branch or projector changes $P$ and possibly
$\Adm$.  Those changes need not preserve $\Image$.

\subsection{Motif coordinates}

For a degree-three embedded junction, a convenient local coordinate chart is
\[
  u=(\beta,q_\parallel,q_\perp,\eta)\in\R^4.
\]
Here $\beta$ measures bending of a declared dominant channel,
$q_\parallel$ and $q_\perp$ are longitudinal and transverse components of a
candidate third branch, and $\eta$ collects any additional measured
attribute such as taper or imbalance.  Define
\[
  \rho=\sqrt{q_\parallel^2+q_\perp^2}.
\]
When $\rho>0$, the branch angle $\phi$ relative to the dominant channel
satisfies
\[
  \cos\phi=\frac{q_\parallel}{\rho}.
\]
Thus $q_\parallel=0$ means an orthogonal branch in this local chart.

These variables are measurements after $\Gamma$ and $m$.  They do not
create a graph edge.  A branch must already have been detected by the
network-extraction rule before its angle and thickness can be measured.

\section{What failure of reconstruction does not prove}

Projection and reconstruction are easily confused.  A section of
$P:\Adm\to P(\Adm)$ is a map $s$ with $P\circ s=\mathrm{id}$.  Its existence
depends on the regularity category: set-theoretic, measurable, continuous,
smooth, or connection-preserving sections are different claims.  Version 1
did not keep those categories separate.

\begin{proposition}[No section does not determine image geometry]
There is a continuous surjection onto a full product space that has no
continuous global section.
\end{proposition}

\begin{proof}
Consider
\[
  p:\Sone\times[-1,1]\longrightarrow\Sone\times[-1,1],
  \qquad p(z,t)=(z^2,t).
\]
It is onto, so its image is the complete product
$\Sone\times[-1,1]$.  If a continuous section existed, the induced maps on
the fundamental group of the circle factor would satisfy
$p_\ast s_\ast=\mathrm{id}_{\mathbb Z}$.  But $p_\ast$ is multiplication by
$2$, so this would require an integer $k$ with $2k=1$, which is impossible.
\end{proof}

The example has two consequences.  First, lack of a continuous section does
not imply that the image is a proper subset of the target.  Second, it does
not imply coupled coordinate constraints.  A reconstruction obstruction
describes the fibers or topology of a map; image geometry must be determined
separately.

The measurable case requires its own theorem and hypotheses.  One must not
take failure of a continuous section and silently promote it to failure of a
measurable section.  Standard measurable-selection results make such a
promotion particularly unsafe \cite{kechris1995}.

\begin{corollary}[Correct use of no-section results]
A no-section theorem can falsify a claim that explicitly requires a section
of the same regularity on the same domain.  By itself it cannot prove that
$\Image$ is proper, nonfactorizing, low-dimensional, graph-like, or
sprout-selecting.
\end{corollary}

\section{Direct tests for hidden compatibility}

\subsection{Support factorization}

Suppose the motif chart is a product
\[
  \Motif=M_1\times\cdots\times M_k
\]
with coordinate projections $\pi_i$.  There are at least two notions of
factorization.

\begin{definition}[Support factorization]
A subset $\Image\subseteq\Motif$ factorizes by coordinates if
\[
  \Image=\prod_{i=1}^k\pi_i(\Image).
\]
Failure of this equality means that some individually realizable coordinate
values cannot be recombined into a jointly realizable motif.
\end{definition}

\begin{definition}[Probabilistic factorization]
A probability law $\mu$ on $\Motif$ factorizes if
\[
  \mu=\mu_1\otimes\cdots\otimes\mu_k
\]
for its coordinate marginals.  This is a statement about a law, not only its
support.
\end{definition}

\begin{proposition}[Recombination certificate]
If $u,v\in\Image$ and a point $w$ obtained by taking at least one coordinate
from $u$ and another from $v$ is not in $\Image$, then $\Image$ does not
factorize.
\end{proposition}

\begin{proof}
Every coordinate of $w$ belongs to the corresponding projection
$\pi_i(\Image)$.  Hence $w$ belongs to
$\prod_i\pi_i(\Image)$.  Since $w\notin\Image$, equality with that product is
impossible.
\end{proof}

This gives a finite falsification certificate when membership in $\Image$
can be checked.  It also identifies what must be computed: two accepted
motifs and one rejected recombination from the same frozen source and
extraction rule.

\subsection{Constraint and graph forms}

An explicit constraint can describe the image locally.  Let
$h:\Motif\to\R^r$ be continuously differentiable and suppose that the
declared image is contained in $h^{-1}(0)$.

\begin{proposition}[Regular local image constraint]
If $0$ is a regular value of $h$, then $h^{-1}(0)$ is a submanifold of
codimension $r$.  If $\Image$ contains a relatively open subset of this
level set, then the stated equations give its local compatibility geometry.
\end{proposition}

\begin{proof}
This is the regular-value theorem \cite{lee2013}.
\end{proof}

Codimension alone does not identify which variables are coupled.  A direct
and useful sufficient condition is a graph relation.

\begin{proposition}[Graph-of-a-map nonfactorization]
Let $U\subseteq M_1$ and let $f:U\to M_2$ be nonconstant.  Then
\[
  \Image_f=\{(x,f(x)):x\in U\}
\]
does not factorize as $\pi_1(\Image_f)\times\pi_2(\Image_f)$.
\end{proposition}

\begin{proof}
Choose $x_1,x_2$ with $f(x_1)\neq f(x_2)$.  Both
$(x_1,f(x_1))$ and $(x_2,f(x_2))$ lie in $\Image_f$, while the recombination
$(x_1,f(x_2))$ does not.  Apply the recombination certificate.
\end{proof}

In an MTT application, $h$ or $f$ must come from a selected source
calculation, a certified reduced equation, or a declared phenomenological
model.  Writing the symbol $\Image$ does not establish any constraint.

\section{Correlation constraints are not graph edges}

Three different uses of the word ``network'' must be separated.

\begin{enumerate}
\item A \emph{physical or geometric graph} has vertices and edges extracted
by $\Gamma$ from an image, flow field, point cloud, or other effective
state.
\item A \emph{motif support} is a set $\Image\subseteq\Motif$ of allowed
measurements.
\item A \emph{statistical graph} encodes conditional independences of a
probability law.
\end{enumerate}

None determines the others without additional assumptions.  In particular,
a correlated pair of motif coordinates is not a new physical edge.
Conversely, two physical edges can carry independent measured attributes.
Graphical-model edges require a specified law and a Markov convention
\cite{lauritzen1996}.

\begin{proposition}[Support does not determine dependence]
The square $[-1,1]^2$ supports both an independent distribution and a
correlated distribution with a strictly positive density everywhere.
\end{proposition}

\begin{proof}
The uniform density $f_0(x,y)=1/4$ is independent.  For
$0<|\epsilon|<1$, define
\[
  f_\epsilon(x,y)=\frac14(1+\epsilon xy).
\]
It is positive and normalized on the same square.  Its marginals are
uniform, but
\[
  \E_\epsilon[XY]=\frac{\epsilon}{9}\neq0.
\]
Thus identical full support is compatible with different dependence.
\end{proof}

\begin{proposition}[Dependence does not determine a physical graph]
Fix any nonproduct law on a two-coordinate motif space.  The same law can be
attached to different embedded graphs by changing the extraction map
$\Gamma$ while leaving the measured coordinate pair unchanged.
\end{proposition}

\begin{proof}
For example, assign the same joint attribute law to two marked branches in a
three-edge star or to two marked edges in a cycle.  The coordinate law is
unchanged while the physical edge sets differ.
\end{proof}

The corrected chain is therefore
\[
  \text{source}
  \xrightarrow{P}
  \text{effective state}
  \xrightarrow{\Gamma}
  \text{physical graph}
  \xrightarrow{m}
  \text{motif coordinates}.
\]
A statistical graph may be inferred only after a law on those coordinates
has also been supplied.

\section{How a local reserve model can arise}

\subsection{Full and reduced quadratic forms}

Near a selected reference configuration, a twice differentiable deficit or
action can be expanded in visible motif coordinates $u$ and hidden response
coordinates $z$:
\[
  D(u,z)
  =
  D(0,0)
  +\frac12
  \begin{pmatrix}u\\ z\end{pmatrix}^{\!T}
  \begin{pmatrix}A&B\\ B^T&D\end{pmatrix}
  \begin{pmatrix}u\\ z\end{pmatrix}
  +R_3(u,z).
\]
The notation $D$ for both the hidden block and the full deficit is avoided
below by writing the hidden block as $H_{zz}$.

\begin{theorem}[Schur-complement reduction]
Let
\[
  H=
  \begin{pmatrix}H_{uu}&H_{uz}\\H_{zu}&H_{zz}\end{pmatrix}
\]
be symmetric with $H_{zz}>0$.  Minimizing the quadratic form over $z$ gives
\[
  \inf_z\frac12
  \begin{pmatrix}u\\z\end{pmatrix}^{\!T}
  H
  \begin{pmatrix}u\\z\end{pmatrix}
  =
  \frac12u^TH_{\mathrm{eff}}u,
\qquad
  H_{\mathrm{eff}}
  =
  H_{uu}-H_{uz}H_{zz}^{-1}H_{zu}.
\]
If $H>0$, then $H_{\mathrm{eff}}>0$.
\end{theorem}

\begin{proof}
Complete the square using
$z_\ast=-H_{zz}^{-1}H_{zu}u$.  Positivity of the Schur complement follows
from positivity of $H$ \cite{hornjohnson2013}.
\end{proof}

This theorem explains how a local visible Hessian can be computed from a
larger source Hessian.  It does not guarantee any particular entries.  In a
physical MTT network model, $H$ must come from the same selected source that
defines $P$ and $\Adm$, and the remainder $R_3$ needs a certified bound.

\subsection{The anisotropic sprout chart}

Take
\[
  u=(\beta,q_\parallel,q_\perp).
\]
The simplest positive quadratic model that penalizes bending, total branch
size, and longitudinal alignment is
\[
  Q(u)
  =
  a\beta^2
  b(q_\parallel^2+q_\perp^2)
  cq_\parallel^2,
  \qquad a,b,c>0.
\]
The coefficient $b$ is the isotropic branch-size cost.  The coefficient $c$
is an additional longitudinal penalty.  Such a diagonal form may result
from symmetry or from diagonalizing a positive local Hessian, but the
identification of the eigenvectors with physical bend and branch directions
is an extra source statement.

For a reserve $C>0$, define the local feasible image
\[
  \Image_C^{(2)}=\{u:Q(u)\leq C\}.
\]
This is an ellipsoid in local motif coordinates.  Its geometry, rather than
projection noninvertibility, produces the following result.

\section{The conditional sprout theorem}

\begin{theorem}[Quadratic reserve bound]
Let $a,b,c,C>0$ and let $u\in\Image_C^{(2)}$.  Then
\[
  \rho\leq\sqrt{\frac{C}{b}}.
\]
Equality holds only when
\[
  \beta=0,\qquad q_\parallel=0,
\]
so a saturating branch leaves the dominant channel straight and is
orthogonal in the declared local chart.

More generally, if $0<\lambda\leq1$ and
\[
  \rho\geq\lambda\sqrt{\frac{C}{b}},
\]
then
\[
  |\beta|
  \leq
  \sqrt{\frac{C(1-\lambda^2)}{a}},
\qquad
  |q_\parallel|
  \leq
  \sqrt{\frac{C(1-\lambda^2)}{c}}.
\]
For $\rho>0$ this gives
\[
  \cos^2\phi
  \leq
  \frac{b(1-\lambda^2)}{c\lambda^2}.
\]
\end{theorem}

\begin{proof}
Since
\[
  C\geq Q
  =a\beta^2+b\rho^2+cq_\parallel^2
  \geq b\rho^2,
\]
the first bound follows.  Equality requires both nonnegative residual terms
to vanish.  Under the lower bound on $\rho$,
\[
  a\beta^2+cq_\parallel^2
  \leq C-b\rho^2
  \leq C(1-\lambda^2).
\]
Bounding each term separately proves the next two inequalities.  Divide the
$q_\parallel$ bound by
$\rho^2\geq\lambda^2C/b$ to obtain the angle estimate.
\end{proof}

The theorem says less, and more precisely, than Version 1.  It does not say
that a branch must appear.  It says that if a branch in this chart uses a
large fraction of the available quadratic reserve, then it must be thin on
the absolute scale $\sqrt C$, nearly transverse, and accompanied by little
dominant-channel bending.

\subsection{Controlled nonlinear remainder}

A physical deficit is rarely exactly quadratic.  Let
\[
  \widetilde Q(u)=Q(u)+R(u).
\]

\begin{theorem}[Relative-error stability]
Suppose $|R(u)|\leq\epsilon Q(u)$ on the declared chart, with
$0\leq\epsilon<1$.  If $\widetilde Q(u)\leq C$, then
\[
  \rho
  \leq
  \sqrt{\frac{C}{(1-\epsilon)b}}.
\]
\end{theorem}

\begin{proof}
The remainder bound gives
$\widetilde Q\geq(1-\epsilon)Q\geq(1-\epsilon)b\rho^2$.
\end{proof}

An additive error $|R|\leq\delta$ instead yields
\[
  \rho\leq\sqrt{\frac{C+\delta}{b}}.
\]
Both forms expose the verification obligation.  One must state the chart,
norm, remainder bound, and source of the Hessian coefficients.

\subsection{Why the exponent is not universal}

\begin{proposition}[Order of contact determines the scaling exponent]
Suppose for small $\rho$ the branch deficit satisfies
\[
  k_-\rho^p\leq D(\rho)\leq k_+\rho^p,
  \qquad p>0,\quad k_\pm>0.
\]
Then the maximal feasible branch scale obeys
\[
  \left(\frac{C}{k_+}\right)^{1/p}
  \leq \rho_{\max}(C)
  \leq
  \left(\frac{C}{k_-}\right)^{1/p}.
\]
\end{proposition}

\begin{proof}
Apply the two inequalities to the condition $D(\rho)\leq C$ and to a trial
value satisfying the upper estimate.
\end{proof}

Thus the exponent $1/2$ is evidence for a nondegenerate quadratic leading
term.  A quartic flat direction gives $C^{1/4}$; a nonsmooth linear cost
gives $C$.  Projection alone selects none of these.

\section{Feasible sets do not supply probabilities}

Version 1 inferred that sprout motifs become statistically dominant near the
boundary.  A feasible set cannot establish that conclusion without a law.

\begin{proposition}[No universal prevalence from support]
Let $\Image_C^{(2)}$ contain both the origin and at least one nonzero
near-saturating sprout.  There are probability measures supported on the
same set for which the sprout probability is $0$, $1$, or any prescribed
number in $[0,1]$.
\end{proposition}

\begin{proof}
Use a Dirac measure at the origin, a Dirac measure at a declared sprout, or
their convex mixtures.
\end{proof}

A statistical theorem therefore needs one of the following:

\begin{itemize}
\item a source probability measure $\nu$ and its pushforward $F_\ast\nu$;
\item a stochastic growth process with a stopping rule;
\item a deterministic ensemble and a declared sampling measure;
\item or a Gibbs, maximum-entropy, or other law justified independently.
\end{itemize}

Even then, conditioning must be specified.  A concentration statement can
be valid under a chosen law while failing under another law with the same
support.  The current MTT network program has not selected such a law.

\section{Boundary language and capacity}

MTT coherence capacity is best treated as a sourced reserve vector or a
declared normalized margin class, not as a universal scalar force.  A local
network model may choose a scalar reserve
\[
  C=\min_i r_i
\]
or another monotone functional of certified reserve rows
$r=(r_1,\ldots,r_n)$.  That scalar is useful for a bottleneck estimate but
does not reconstruct the rows or their Hessian.

For this paper, $C$ has only the following role:

\begin{quote}
$C$ is the certified amount by which the declared local deficit may increase
before leaving the selected admissible chart.
\end{quote}

It is not automatically energy, curvature, crowding, entropy, stress, or
biological fitness.  A proposed observable proxy $\widehat C$ must be
calibrated against the reserve it is intended to estimate.  Monotonicity
between a proxy and a geometric feature cannot be used both to define
capacity and to confirm the prediction without circularity.

\section{A valid empirical test contract}

The corrected model can be tested, but only with a predeclared data-use
contract.

\subsection{Training-stage obligations}

\begin{enumerate}
\item Fix the network-extraction algorithm $\Gamma$, including resolution,
pruning, junction, and thickness rules.
\item Fix the motif map $m$ and the local coordinate chart
$(\beta,q_\parallel,q_\perp,\eta)$.
\item State how a candidate reserve or reserve proxy is obtained.
\item Estimate or derive the local Hessian and verify $a,b,c>0$ on the
training domain.
\item Bound the nonlinear remainder and freeze the accepted chart.
\end{enumerate}

\subsection{Held-out obligations}

On data not used to choose the extraction, proxy, chart, or coefficients,
test:

\begin{enumerate}
\item whether observed accepted motifs lie inside the certified feasible
region up to measurement error;
\item whether the upper envelope of $\rho$ follows the predicted reserve
scaling;
\item whether high-reserve-fraction branches obey the bend and angle bounds;
\item whether the same coefficients transfer across the preregistered
conditions claimed by the model.
\end{enumerate}

Failure of a bound with verified premises falsifies the local reserve model
on that domain.  Failure of an uncalibrated proxy does not falsify MTT as a
whole.  Conversely, a fitted square-root envelope is not a source theorem:
many positive quadratic models share it.

\subsection{What would count as strong evidence}

The strongest version would derive $F$, $H_{\mathrm{eff}}$, and the reserve
rows from one selected source before the network measurements are used.
The theorem would then predict coefficients and transfer behavior
out-of-sample.  A weaker but still useful phenomenological test can fit
$a,b,c$ on one condition and validate them on another.  The two evidence
tiers must not be merged.

\section{Relation to established network models}

This framework does not show that optimization is unnecessary in real
networks.  Established models demonstrate that optimization, damage
tolerance, fluctuating loads, local adaptation, and growth can alter network
topology and hierarchy \cite{katifori2010,corson2010,ronellenfitsch2016}.
Those results are mechanisms with their own state variables and objectives.

The MTT reserve model is different.  It asks whether a selected source and
projection induce a constrained motif image whose local Hessian has a
specific anisotropy.  If so, the conditional sprout theorem follows without
minimizing a global network functional.  But the source of the local
quadratic form still needs explanation.  In some systems it may arise from
elasticity, transport, growth, or an optimization principle.  MTT does not
erase those domain dynamics.

Nor does the framework establish a string-theory description.  A minimal
surface or worldsheet model can be compared only after an explicit map
identifies its fields, boundary conditions, and observables with the same
network construction.  Similar-looking geometry is not an equivalence
theorem.

\clearpage
\section{Current MTT status}

Table~\ref{tab:status} prevents the conditional theorem from being read as a
selected physical result.

\begin{table}[ht]
\centering
\small
\begin{tabularx}{\textwidth}{@{}p{0.25\textwidth}p{0.20\textwidth}X@{}}
\toprule
Object & Current tier & What is established or missing\\
\midrule
Typed chain $F=m\circ\Gamma\circ P$
& framework
& A valid way to state a network encoding; no physical network instance is
selected by the notation.\\
Image tests
& exact mathematics
& Recombination, graph-of-a-map, and regular-value tests are rigorous once
membership or constraints are supplied.\\
No-section implication
& withdrawn
& No-section alone does not imply a proper or nonfactorizing image.\\
Quadratic sprout bound
& conditional exact
& Exact for the declared $Q$ and controlled under the stated remainder
bound.\\
Physical network source
& open
& No selected MTT operator or action currently emits $\Gamma$, the motif
chart, and the local network Hessian from one source.\\
Coefficients $a,b,c$
& open
& Not derived for a biological, physical, or engineered network.\\
Capacity calibration
& open
& Curvature, crowding, and taper are not accepted universal capacity
proxies.\\
Prevalence or growth law
& open
& A feasible image does not select a probability measure or branch-creation
dynamics.\\
Cross-domain universality
& not established
& Similar motifs across domains do not establish a common source.\\
\bottomrule
\end{tabularx}
\caption{The corrected claim ledger.}
\label{tab:status}
\end{table}

The model is therefore a conditional encoding at the representation and
local reconstruction level.  It is not a derivation of observed network
statistics from current MTT geometry.

\clearpage
\section{Version 2 changes and reasons}

\begin{table}[ht]
\centering
\small
\begin{tabularx}{\textwidth}{@{}p{0.30\textwidth}p{0.28\textwidth}X@{}}
\toprule
Version 1 statement & Version 2 decision & Reason\\
\midrule
No global reconstruction forces hidden motif constraints
& Withdraw
& A surjective map without a continuous section can have a full product
image.\\
Local observables cannot be freely combined
& Replace by a direct image test
& Nonfactorization requires a failed recombination or an explicit coupled
constraint.\\
Correlation constraints generate network geometry
& Separate the objects
& Physical edges come from $\Gamma$; statistical edges require a law and
conditional-independence convention.\\
Sprouting is the only admissible boundary deformation
& Narrow to a local quadratic theorem
& The conclusion needs a positive anisotropic Hessian and a controlled
remainder.\\
$\rho_{\rm th}\sim\sqrt C$ is universal
& Make conditional
& The exponent is $1/p$ when the first nonzero branch cost has order $p$.\\
Sprouts dominate boundary ensembles
& Withdraw without a measure
& The same support admits probability laws with any sprout prevalence.\\
Capacity proxies test the theory directly
& Require calibration and held-out data
& An observable proxy is not automatically the sourced reserve.\\
Optimization and string descriptions are unnecessary
& Remove as a general claim
& Other mechanisms can generate network structure; comparison requires an
explicit common map.\\
\bottomrule
\end{tabularx}
\caption{Claim-by-claim revision ledger.}
\end{table}
\clearpage

\section{Conclusion}

The central correction is simple: information loss is not geometry.
Projection may make reconstruction nonunique or obstruct regular sections,
but the observable network constraints live in the actual image
$\Image=F(\Adm)$.  They must be computed or specified there.

Once that is done, MTT can support a clean local theorem.  A selected
positive Hessian with an extra longitudinal penalty defines an ellipsoidal
reserve region.  The maximal branch scale is then proportional to
$\sqrt C$, and branches near that scale are forced to be approximately
orthogonal while the dominant channel remains nearly straight.  A relative
remainder certificate preserves the conclusion with an explicit loss.

This is a useful result, but it is conditional.  It neither creates the
branch nor predicts its frequency.  It does not choose graph edges, source
coefficients, a biological capacity proxy, or a cross-domain growth law.
Those are precisely the objects that a future selected MTT network model
must emit.  By separating them, Version 2 turns a broad analogy into an
auditable research program with exact mathematical tests and clear failure
conditions.

\bibliographystyle{plain}
\bibliography{main}

\end{document}
